\begin{qfdefinition}[Natural gas for macro traders: producers, shale, flaring, LNG and the Henry Hub futures contract as Gliner describes them in 2013]\label{def:natural-gas-market}
\qffield{Statement} Gliner calls natural gas ``the cleanest and safest fossil fuel'', producing ``30 to 50 percent less carbon dioxide than gasoline or coal''; it is used in power generation, heating, transport and to make ammonia for fertiliser \citep[chap.~10, sec.\ ``Natural Gas'', PDF pp.~313--316]{s7-gliner}. Shale gas, produced with horizontal drilling and hydraulic fracturing, is expected to supply 50 percent of U.S.\ gas by 2020 (p.~314). Flaring burns unwanted gas, abundant in the United States where it is produced alongside oil even when there is no demand (p.~315). Cooled to $-260^\circ$F, gas becomes LNG, taking ``600 times less space'' (p.~315). The United States and Russia account for 38 percent of world production (Table~10.5, p.~316).

Gas is priced in dollars per million British thermal units (MMBtu), roughly one thousand cubic feet (p.~317); the NYMEX contract (NG) is for delivery at Henry Hub, Louisiana; the front of the curve is ``predominantly a speculators' market'' and the back a producers' hedging market; the UNG ETF tracks the contract, with contango and rolls a caution (pp.~317--318).

\qffield{Notation} MMBtu million British thermal units; Mcf thousand cubic feet; LNG liquefied natural gas.

\qffield{Conditions} Figures are Gliner's for 2012--2013. Emissions comparisons depend on whether methane leakage is counted. Gliner gives the contract as ``10,000 million British thermal units'' (10,000 MMBtu, p.~317), as of 2013, which the current chapter confirms (below). \emph{Delivery, as amended in 2016}. NYMEX Rule~220104 of the Henry Hub Natural Gas futures (NG), as certified to the CFTC on 11 July 2016 with effect from 26 July 2016, provides that ``Delivery shall be made free-on-board (`F.O.B.') at the buyer's interconnection point at the Henry Hub'', is complete when gas passes that point ``via physical flow or displacement'', and recognises the Intra Hub Transfer service as an acceptable interconnection point; the filing describes the hub as operated by Sabine Pipe Line \citep[Submission 16-080, pp.~1--3 and Appendix~A]{nymex-16-080}.

\emph{Current chapter} (official PDF as of October 2026; no amendment date in the text, and its page footer reads ``\copyright\ Copyright 2009''). The trading unit is 10{,}000 MMBtu with a delivery tolerance of 2 percent either way; the minimum fluctuation is \$0.001 per MMBtu (\$10 a contract), with \$0.00025 for inter-commodity spreads on Globex; trading ends on the third business day before the first day of the delivery month; and Rule~220104 reads as certified in 2016, with Sabine Pipe Line LLC as the hub operator \citep[Rules 220101--220104]{nymex-ch220-2026}. Whether Chapter~220 was amended after 2016 in respects not covered here cannot be told from the undated text.

Filings (not an exhaustive list): NYMEX Submission 17-048 of 23 February 2017 certified ``administrative, non-substantive amendments'' to ten natural gas contracts, Chapter~220 among them, effective for trade date 13 March 2017, which ``will, in no manner, impact the structure of the Contracts'' \citep[pp.~1--2]{nymex-17-048}; Submission 18-465 of 15 November 2018 amended Rule~589's special price-fluctuation limits for NG and seven associated contracts \citep[cover sheet]{nymex-18-465}, a Chapter~5 rule, not the Chapter~220 text. Two 2019 filings are relevant. Submission 19-158 expanded trading at settlement (TAS) for Henry Hub futures from ``the spot (except on the last trading day) and second through seventh contract months'' to the spot and ``second through twelve contract months'', for trade date 20 May 2019 \citep[pp.~1--2]{nymex-19-158}.

Submission 19-198 standardised the daily settlement procedure documents of six Henry Hub contracts, Chapter~220 included, from trade date 4 September 2019 \citep[pp.~1--2]{nymex-19-198}. Both are Chapter~5 or procedure changes, not amendments to the Chapter~220 text. \emph{After 2019}. The CFTC's listing of NYMEX product terms-and-conditions filings from September 2019 to October 2026 holds 281 filings. Submission 20-166 of 20 March 2020 raised the special price fluctuation limits of Rule~589 for NG and five associated contracts ``from 7\% to 15\% in light of recent market volatility'', effective for trade date 6 April 2020 \citep[pp.~1--2]{nymex-20-166}. That is again a Chapter~5 change. The other natural-gas filings in the period concern options, Platts and TTF contracts, or other look-alikes.

None amends the Chapter~220 text after 2019. The listing covers product terms only, not organisation-rule filings, and the current chapter is undated, so an unlisted change cannot be excluded.

\qffield{Practice} Henry Hub is one hub; regional hubs and LNG-linked prices can diverge from it. \emph{U.S. LNG exports since Gliner}: EIA's annual series of liquefied U.S. natural gas exports, in million cubic feet, runs 2{,}924 in 2013 and 16{,}255 in 2014, 186{,}841 in 2016, 1{,}819{,}547 in 2019 and 5{,}508{,}984 in 2025, about 15.1 billion cubic feet a day \citep[series N9133US2]{eia-lng-exports-2026}, by the release of 30 September 2026. So the market Gliner described in 2013, with LNG essentially absent from U.S. exports, has become one in which LNG exports link Henry Hub to world gas prices (Interpretation; the link itself is not measured here).

\qffield{Pitfalls} Attributing U.S.\ flaring to regulatory requirements is misleading; Gliner attributes it to gas produced where there is no demand or infrastructure. That shale turned the U.S.\ from importer into exporter goes beyond Gliner's 2013 account.

\qffield{See also} \hyperref[def:crude-oil-markets]{crude oil markets}; \hyperref[def:contango-backwardation-stocks]{contango and backwardation}.

\qffield{Further reading} Hub formation and gas--oil price links.

\end{qfdefinition}
